% GENERATED FILE. Edit canonical lessons, then run npm run generate:books.
% canonical-source-hash: fnv1a64:1be6e8c791e84590
% canonical-lessons: MR-C450-man, MR-C450-jiv, MR-C450-sukh, MR-C450-ascarya, MR-C450-helavne

\chapter{Mind and Heart}
\label{ch:mr-mind-and-heart}
\hlchaptermodality{450}

\begin{chapteropening}
\textbf{By the end of this chapter:} \emph{I can talk about the mind, life, happiness and surprise.}

Complete the last lesson of chapter 450: I can talk about the mind, life, happiness and surprise.
\end{chapteropening}

\section[\texorpdfstring{\mr{मन} (man)}{मन (man)}]{\mr{मन} (man) — the mind; the heart, as the seat of feeling}

\label{lesson:MR-C450-man}



\begin{quote}
Before the new one: say the Marathi for dance, then the Marathi for a stage.
\end{quote}

\subsection*{You'll want to know: \mr{मन}}
\textbf{\mr{मन}} — \emph{man} — \textquotedblleft{}the mind; the heart, as the seat of feeling\textquotedblright{}.

\begin{quote}
\mr{आज} \mr{माझं} \mr{मन} \mr{शांत} \mr{आहे}.
\end{quote}

\begin{quote}
\emph{āj mājhaṁ man śānt āhe.}
\end{quote}

\begin{quote}
My mind is calm today.
\end{quote}

Neuter.

Neuter: \mr{माझं} \mr{मन}. \mr{मनात}, in the mind; \mr{मनापासून}, from the heart.

\begin{cousinweb}
\textbf{\mr{मन}} (\emph{man}) is a tatsama, taken from Sanskrit \emph{manas}, \textquotedblleft{}mind\textquotedblright{} (CDIAL 9822); Marathi keeps the stem and drops the final \emph{-s}.
\end{cousinweb}

\begin{grammarlens}[title={What you've built}]
One more word for describing what you lived and felt.
\end{grammarlens}

\subsection*{Your turn}
\begin{itemize}
\raggedright
  \item \emph{Say it:} \emph{man}
  \item \emph{Say it:} \emph{man}, once more
  \item \emph{Say it:} \emph{āj mājhaṁ man śānt āhe.} — My mind is calm today.
  \item \emph{Recall:} say the Marathi for dance, then the Marathi for a stage, then say \emph{man} again
\end{itemize}

\begin{tcolorbox}[breakable,colback=teal!4,colframe=teal!35!black,title={Before you move on}]
What does \mr{मन} mean? (\textquotedblleft{}The mind; the heart, as the seat of feeling\textquotedblright{}.) Say it once more.
\end{tcolorbox}
\section[\texorpdfstring{\mr{जीव} (jīv)}{जीव (jīv)}]{\mr{जीव} (jīv) — life; a living being}

\label{lesson:MR-C450-jiv}



\begin{quote}
Before the new one: say the Marathi for a stage, then the Marathi for the mind.
\end{quote}

\subsection*{You'll want to know: \mr{जीव}}
\textbf{\mr{जीव}} — \emph{jīv} — \textquotedblleft{}life; a living being\textquotedblright{}.

\begin{quote}
\mr{त्या} \mr{लहान} \mr{पक्ष्यात} \mr{अजून} \mr{जीव} \mr{होता}.
\end{quote}

\begin{quote}
\emph{tyā lahān pakṣyāt ajūn jīv hotā.}
\end{quote}

\begin{quote}
That little bird still had life in it.
\end{quote}

Masculine.

Masculine. \mr{जीव} \mr{लावणे}, to grow fond of someone; \mr{जिवंत}, alive.

\begin{cousinweb}
\textbf{\mr{जीव}} (\emph{jīv}) is a tatsama, Sanskrit \emph{jīva}, \textquotedblleft{}alive, a living being\textquotedblright{}; CDIAL lists Marathi \emph{jīv}, \textquotedblleft{}living being, energy\textquotedblright{}, under it (CDIAL 5239).
\end{cousinweb}

\begin{grammarlens}[title={What you've built}]
One more word for describing what you lived and felt.
\end{grammarlens}

\subsection*{Your turn}
\begin{itemize}
\raggedright
  \item \emph{Say it:} \emph{jīv}
  \item \emph{Say it:} \emph{jīv}, once more
  \item \emph{Say it:} \emph{tyā lahān pakṣyāt ajūn jīv hotā.} — That little bird still had life in it.
  \item \emph{Recall:} say the Marathi for a stage, then the Marathi for the mind, then say \emph{jīv} again
\end{itemize}

\begin{tcolorbox}[breakable,colback=teal!4,colframe=teal!35!black,title={Before you move on}]
What does \mr{जीव} mean? (\textquotedblleft{}Life; a living being\textquotedblright{}.) Say it once more.
\end{tcolorbox}
\section[\texorpdfstring{\mr{सुख} (sukh)}{सुख (sukh)}]{\mr{सुख} (sukh) — happiness; comfort}

\label{lesson:MR-C450-sukh}



\begin{quote}
Before the new one: say the Marathi for the mind, then the Marathi for life.
\end{quote}

\subsection*{You'll want to know: \mr{सुख}}
\textbf{\mr{सुख}} — \emph{sukh} — \textquotedblleft{}happiness; comfort\textquotedblright{}.

\begin{quote}
\mr{सुख} \mr{पैशात} \mr{नसतं}, \mr{मनात} \mr{असतं}.
\end{quote}

\begin{quote}
\emph{sukh paiśāt nastaṁ, manāt astaṁ.}
\end{quote}

\begin{quote}
Happiness is not in money, it is in the mind.
\end{quote}

Neuter.

Neuter, the opposite of \mr{दुःख}. \mr{सुखी}, happy and content.

\begin{cousinweb}
\textbf{\mr{सुख}} (\emph{sukh}) is a tatsama, taken from Sanskrit \emph{sukha}, \textquotedblleft{}comfort, happiness\textquotedblright{} (CDIAL 13451).
\end{cousinweb}

\begin{grammarlens}[title={What you've built}]
One more word for describing what you lived and felt.
\end{grammarlens}

\subsection*{Your turn}
\begin{itemize}
\raggedright
  \item \emph{Say it:} \emph{sukh}
  \item \emph{Say it:} \emph{sukh}, once more
  \item \emph{Say it:} \emph{sukh paiśāt nastaṁ, manāt astaṁ.} — Happiness is not in money, it is in the mind.
  \item \emph{Recall:} say the Marathi for the mind, then the Marathi for life, then say \emph{sukh} again
\end{itemize}

\begin{tcolorbox}[breakable,colback=teal!4,colframe=teal!35!black,title={Before you move on}]
What does \mr{सुख} mean? (\textquotedblleft{}Happiness; comfort\textquotedblright{}.) Say it once more.
\end{tcolorbox}
\section[\texorpdfstring{\mr{आश्चर्य} (āścarya)}{आश्चर्य (āścarya)}]{\mr{आश्चर्य} (āścarya) — surprise; astonishment}

\label{lesson:MR-C450-ascarya}



\begin{quote}
Before the new one: say the Marathi for life, then the Marathi for happiness.
\end{quote}

\subsection*{You'll want to know: \mr{आश्चर्य}}
\textbf{\mr{आश्चर्य}} — \emph{āścarya} — \textquotedblleft{}surprise; astonishment\textquotedblright{}.

\begin{quote}
\mr{मीराला} \mr{पाहून} \mr{मला} \mr{खूप} \mr{आश्चर्य} \mr{वाटलं}.
\end{quote}

\begin{quote}
\emph{mīrālā pāhūn malā khūp āścarya vāṭlaṁ.}
\end{quote}

\begin{quote}
I was very surprised to see Meera.
\end{quote}

Neuter.

\mr{आश्चर्य} \mr{वाटणे}, to be surprised, with the dative: \mr{मला} \mr{आश्चर्य} \mr{वाटलं}. \mr{आश्चर्यकारक}, amazing.

\begin{cousinweb}
\textbf{\mr{आश्चर्य}} (\emph{āścarya}) is a tatsama, taken from Sanskrit \emph{āścarya}, \textquotedblleft{}wonderful\textquotedblright{} (CDIAL 1464); in Marathi it is a neuter noun, \textquotedblleft{}surprise\textquotedblright{}.
\end{cousinweb}

\begin{grammarlens}[title={What you've built}]
One more word for describing what you lived and felt.
\end{grammarlens}

\subsection*{Your turn}
\begin{itemize}
\raggedright
  \item \emph{Say it:} \emph{āścarya}
  \item \emph{Say it:} \emph{āścarya}, once more
  \item \emph{Say it:} \emph{mīrālā pāhūn malā khūp āścarya vāṭlaṁ.} — I was very surprised to see Meera.
  \item \emph{Recall:} say the Marathi for life, then the Marathi for happiness, then say \emph{āścarya} again
\end{itemize}

\begin{tcolorbox}[breakable,colback=teal!4,colframe=teal!35!black,title={Before you move on}]
What does \mr{आश्चर्य} mean? (\textquotedblleft{}Surprise; astonishment\textquotedblright{}.) Say it once more.
\end{tcolorbox}
\section[\texorpdfstring{\mr{हेलावणे} (helāvṇe)}{हेलावणे (helāvṇe)}]{\mr{हेलावणे} (helāvṇe) — to be deeply moved}

\label{lesson:MR-C450-helavne}



\begin{quote}
Before the new one: say the Marathi for happiness, then the Marathi for surprise.
\end{quote}

\subsection*{You'll want to know: \mr{हेलावणे}}
\textbf{\mr{हेलावणे}} — \emph{helāvṇe} — \textquotedblleft{}to be deeply moved\textquotedblright{}.

\begin{quote}
\mr{आईचं} \mr{पत्र} \mr{वाचून} \mr{तो} \mr{हेलावला}.
\end{quote}

\begin{quote}
\emph{āīcaṁ patra vācūn to helāvlā.}
\end{quote}

\begin{quote}
Reading his mother's letter, he was deeply moved.
\end{quote}

Past (he …): \mr{तो} \mr{हेलावला}, \emph{to helāvlā}.

Often with \mr{मन}: \mr{तिचं} \mr{मन} \mr{हेलावलं}, her heart was stirred.

\begin{cousinweb}
\textbf{\mr{हेलावणे}} (\emph{helāvṇe}) is a Marathi verb whose origin is not settled.
\end{cousinweb}

\begin{grammarlens}[title={What you've built}]
One more word for describing what you lived and felt.
\end{grammarlens}

\subsection*{Your turn}
\begin{itemize}
\raggedright
  \item \emph{Say it:} \emph{helāvṇe}
  \item \emph{Say it:} \emph{helāvṇe}, once more
  \item \emph{Say it:} \emph{āīcaṁ patra vācūn to helāvlā.} — Reading his mother's letter, he was deeply moved.
  \item \emph{Recall:} say the Marathi for happiness, then the Marathi for surprise, then say \emph{helāvṇe} again
\end{itemize}

\begin{tcolorbox}[breakable,colback=teal!4,colframe=teal!35!black,title={Before you move on}]
What does \mr{हेलावणे} mean? (\textquotedblleft{}To be deeply moved\textquotedblright{}.) Say it once more.
\end{tcolorbox}
